\documentclass[11pt]{article}
\usepackage[vietnamese]{babel}
\usepackage[utf8]{inputenc}
\usepackage[T5]{fontenc}
\usepackage{amsmath}
\usepackage{amsfonts}
\usepackage{amssymb}
\usepackage[version=4]{mhchem}
\usepackage{stmaryrd}
\usepackage{graphicx}
\usepackage[export]{adjustbox}
\graphicspath{ {./images/} }
\usepackage{bbold}

\begin{document}
BỘ GIÁO DỤC VÀ ĐÀO TẠO

ĐỀ THI CHÍNH THÚC

KỲ THI THPT QUỐC GIA NĂM 2019\\
Bài thi: TOÁN HỌC\\
Thời gian làm bài: 90 phút, không kể thời gian phát đề

MÃ ĐỀ 101\\
Câu 1. Trong không gian $O x y z$, cho mặt phẳng $(P): x+2 y+3 z-1=0$. Vectơ nào dưới đây là một vectơ pháp tuyến của $(P)$ ?\\
A. $\overrightarrow{n_{3}}=(1 ; 2 ;-1)$.\\
B. $\overrightarrow{n_{4}}=(1 ; 2 ; 3)$.\\
C. $\overrightarrow{n_{1}}=(1 ; 3 ;-1)$.\\
D. $\overrightarrow{n_{2}}=(2 ; 3 ;-1)$.

Câu 2. Với $a$ là số thực dương tùy, $\log _{5} a^{2}$ bằng\\
A. $2 \log _{5} a$.\\
B. $2+\log _{5} a$.\\
C. $\frac{1}{2}+\log _{5} a$.\\
D. $\frac{1}{2} \log _{5} a$.

Câu 3. Cho hàm số $f(x)$ có bảng biến thiên như sau:

\begin{center}
\begin{tabular}{r|rrcccccc}
$x$ & $-\infty$ & -2 & 0 &  & 2 & $+\infty$ &  &  \\
$f^{\prime}(x)$ &  & - & 0 & + & 0 & - & 0 & + \\
\hline
$f(x)$ & $+\infty$ &  &  &  & 3 &  &  & $+\infty$ \\
\hline
\end{tabular}
\end{center}

Hàm số đã cho nghịch biến trên khoảng nào dưới đây?\\
A. (-2;0).\\
B. $(2 ;+\infty)$.\\
C. ( $0 ; 2$ ).\\
D. $(0 ;+\infty)$.

Câu 4. Nghiệm phương trình $3^{2 x-1}=27$ là\\
A. $x=5$.\\
B. $x=1$.\\
C. $x=2$.\\
D. $x=4$.

Câu 5. Cho cấp số cộng $\left(u_{n}\right)$ với $u_{1}=3$ và $u_{2}=9$. Công sai của cấp số cộng đã cho bằng\\
A. -6 .\\
B. 3 .\\
C. 12 .\\
D. 6 .

Câu 6. Đồ thị của hàm số nào dưới đây có dạng như đường cong hình vẽ bên\\
\includegraphics[max width=\textwidth, alt={}, center]{6a73def7-dd88-4569-925a-9b949b4179bd-1_403_343_1665_884}\\
A. $y=x^{3}-3 x^{2}+3$.\\
B. $y=-x^{3}+3 x^{2}+3$.\\
C. $y=x^{4}-2 x^{2}+3$.\\
D. $y=-x^{4}+2 x^{2}+3$.

Câu 7. Trong không gian $O x y z$, cho đường thẳng $d: \frac{x-2}{-1}=\frac{y-1}{2}=\frac{z+3}{1}$. Vectơ nào dưới đây là một vectơ chỉ phương của $d$ ?\\
A. $\overrightarrow{u_{2}}=(2 ; 1 ; 1)$. .\\
B. $\overrightarrow{u_{4}}=(1 ; 2 ;-3)$..\\
C. $\overrightarrow{u_{3}}=(-1 ; 2 ; 1) .$.\\
D. $\overrightarrow{u_{1}}=(2 ; 1 ;-3) .$.

Câu 8. Thể tích của khối nón có chiều cao $h$ và bán kính $r$ là\\
A. $\frac{1}{3} \pi r^{2} h$. .\\
B. $\pi r^{2} h$. .\\
C. $\frac{4}{3} \pi r^{2} h .$.\\
D. $2 \pi r^{2} h .$.

Câu 9. Số cách chọn 2 học sinh từ 7 học sinh là\\
A. $2^{7}$.\\
B. $A_{7}^{2}$.\\
C. $C_{7}^{2}$.\\
D. $7^{2}$.

Câu 10. Trong không gian $O x y z$, hình chiếu vuông góc của điểm $M(2 ; 1 ;-1)$ trên trục $O z$ có tọa độ là\\
A. $(2 ; 1 ; 0)$.\\
B. $(0 ; 0 ;-1)$.\\
C. $(2 ; 0 ; 0)$.\\
D. ( $0 ; 1 ; 0$ ).

Câu 11. Biết $\int_{0}^{1} f(x) d x=-2$ và $\int_{0}^{1} g(x) d x=3$, khi đó $\int_{0}^{1}[f(x)-g(x)] d x$ bằng\\
A. -5. .\\
B. 5 . .\\
C. -1. .\\
D. 1..

Câu 12. Thể tích khối lăng trụ có diện tích đáy $B$ và chiều cao $h$ là\\
A. $3 B h$. .\\
B. Bh..\\
C. $\frac{4}{3} B h$. .\\
D. $\frac{1}{3} B h$.

Câu 13. Số phức liên hợp của số phức $3-4 i$ là\\
A. $-3-4 i$.\\
B. $-3+4 i$.\\
C. $3+4 i$.\\
D. $-4+3 i$.

Câu 14. Cho hàm số $f(x)$ có bảng biến thiên như sau:

\begin{center}
\begin{tabular}{r|rrccccc}
$x$ & $-\infty$ &  & -1 &  & 2 &  & $+\infty$ \\
\hline
$f^{\prime}(x)$ &  & - & 0 & + & 0 & - &  \\
\hline
 & $+\infty$ &  &  &  & 1 &  &  \\
\hline
\end{tabular}
\end{center}

Hàm số đã cho đạt cực tiểu tại\\
A. $x=2$.\\
B. $x=1$.\\
C. $x=-1$.\\
D. $x=-3$.

Câu 15. Họ tất cả các nguyên hàm của hàm số $f(x)=2 x+5$ là\\
A. $x^{2}+5 x+C .$.\\
B. $2 x^{2}+5 x+C$. .\\
C. $2 x^{2}+C$.\\
D. $x^{2}+C$..

Câu 16. Cho hàm số $f(x)$ có bảng biến thiên như sau:

\begin{center}
\begin{tabular}{|l|l|l|l|l|l|l|}
\hline
$x$ & $-\infty$ & -2 & 0 &  & 2 & $+\infty$ \\
\hline
$f^{\prime}(x)$ &  & 0 & 0 & + & 0 &  \\
\hline
$f(x)$ &  & 3 & -1 &  & 3 & $-\infty$ \\
\hline
\end{tabular}
\end{center}

Số nghiệm thực của phương trình $2 f(x)-3=0$ là\\
A. 2 .\\
B. 1 .\\
C. 4 .\\
D. 3 .

Câu 17. Cho hình chóp $S . A B C$ có $S A$ vuông góc với mặt phẳng ( $A B C$ ), $S A=2 a$, tam giác $A B C$ vuông tại $B, A B=a \sqrt{3}$ và $B C=a$ (minh họa hình vẽ bên). Góc giữa đường thẳng $S C$ và mặt phẳng ( $A B C$ ) bằng\\
\includegraphics[max width=\textwidth, alt={}, center]{6a73def7-dd88-4569-925a-9b949b4179bd-2_547_574_1841_767}\\
A. $90^{\circ}$.\\
B. $45^{\circ}$.\\
C. $30^{\circ}$.\\
D. $60^{\circ}$.

Câu 18. Gọi $z_{1}, z_{2}$ là hai nghiệm phức phương trình $z^{2}-6 z+10=0$. Giá trị $z_{1}^{2}+z_{2}^{2}$ bằng\\
A. 16 .\\
B. 56 .\\
C. 20 .\\
D. 26 .

Câu 19. Cho hàm số $y=2^{x^{2}-3 x}$ có đạo hàm là\\
A. $(2 x-3) \cdot 2^{x^{2}-3 x} \cdot \ln 2$.\\
B. $2^{x^{2}-3 x} \cdot \ln 2$.\\
C. $(2 x-3) \cdot 2^{x^{2}-3 x}$.\\
D. $\left(x^{2}-3 x\right) \cdot 2^{x^{2}-3 x-1}$.

Câu 20. Giá trị lớn nhất của hàm số $f(x)=x^{3}-3 x+2$ trên đoạn $[-3 ; 3]$ bằng\\
A. -16 .\\
B. 20 .\\
C. 0 .\\
D. 4 .

Câu 21. Trong không gian $O x y z$, cho mặt cầu $(S): x^{2}+y^{2}+z^{2}+2 x-2 z-7=0$. bán kính của mặt cầu đã cho bằng\\
A. $\sqrt{7}$.\\
B. 9 .\\
C. 3 .\\
D. $\sqrt{15}$.

Câu 22. Cho khối lăng trụ đứng $A B C \cdot A^{\prime} B^{\prime} C^{\prime}$ có đáy là tam giác đều cạnh $a$ và $A A^{\prime}=\sqrt{3} a$ (hình minh họa như hình vẽ). Thể tích của lăng trụ đã cho bằng\\
\includegraphics[max width=\textwidth, alt={}, center]{6a73def7-dd88-4569-925a-9b949b4179bd-3_467_456_358_817}\\
A. $\frac{3 a^{3}}{4}$.\\
B. $\frac{3 a^{3}}{2}$.\\
C. $\frac{a^{3}}{4}$.\\
D. $\frac{a^{3}}{2}$.

Câu 23. Cho hàm số $f(x)$ có đạo hàm $f^{\prime}(x)=x(x+2)^{2}, \forall x \in \mathbb{R}$. Số điểm cực trị của hàm số đã cho là\\
A. 0 .\\
B. 3 .\\
C. 2 .\\
D. 1 .

Câu 24. Cho $a$ và $b$ là hai số thực dương thỏa mãn $a^{4} b=16$. Giá trị của $4 \log _{2} a+\log _{2} b$ bằng\\
A. 4 .\\
B. 2 .\\
C. 16 .\\
D. 8 .

Câu 25. Cho hai số phức $z_{1}=1-i$ và $z_{2}=1+2 i$. Trên mặt phẳng toạ độ $O x y$, điểm biểu diễn số phức $3 z_{1}+z_{2}$ có toạ độ là\\
A. $(4 ;-1)$.\\
B. (-1;4).\\
C. (4;1).\\
D. $(1 ; 4)$.

Câu 26. Nghiệm của phương trình $\log _{3}(x+1)+1=\log _{3}(4 x+1)$ là\\
A. $x=3$.\\
B. $x=-3$.\\
C. $x=4$.\\
D. $x=2$.

Câu 27. Một cở sở sản xuất có hai bể nước hình trụ có chiều cao bằng nhau, bán kính đáy lần lượt bằng $1 m$ và $1,2 m$. Chủ cơ sở dự định làm một bể nước mới, hình trụ, có cùng chiều cao và có thể tích bằng tổng thể tích của hai bể nước trên. Bán kính đáy của bể nước dự dịnh làm gần nhất với kết quả nào dưới đây?\\
A. $1,8 m$. .\\
B. $1,4 m$. .\\
C. $2,2 m$. .\\
D. $1,6 m$. .

Câu 28. Cho hàm số $y=f(x)$ có bảng biến thiên như sau:

\begin{center}
\begin{tabular}{c|cc||cccc}
$x$ & $-\infty$ & \multicolumn{2}{c}{0} & 1 & $+\infty$ \\
\hline
$y^{\prime}$ & - & - & 0 & + &  \\
\hline
$y$ & 2 &  & $+\infty$ &  & $+{ }^{+\infty}$ \\
\hline
\end{tabular}
\end{center}

Tổng số tiệm cận đứng và tiệm cận ngang của đồ thị hàm số đã cho là\\
A. 4. .\\
B. 1. .\\
C. 3. .\\
D. 2. .

Câu 29. Cho hàm số $f(x)$ liên tục trên $\mathbb{R}$. Gọi $S$ là diện tích hình phẳng giới hạn bởi các đường $y=f(x), y=0, x=-1$ và $x=4$ (như hình vẽ bên). Mệnh đề nào dưới đây là đúng?\\
\includegraphics[max width=\textwidth, alt={}, center]{6a73def7-dd88-4569-925a-9b949b4179bd-3_387_565_2368_774}\\
A. $S=-\int_{-1}^{1} f(x) d x+\int_{1}^{4} f(x) d x$.\\
B. $S=\int_{-1}^{1} f(x) d x-\int_{1}^{4} f(x) d x$.\\
C. $S=\int_{-1}^{1} f(x) d x+\int_{1}^{4} f(x) d x$.\\
D. $S=-\int_{-1}^{1} f(x) d x-\int_{1}^{4} f(x) d x$.

Câu 30. Trong không gian $O x y z$, cho hai điểm $A(1 ; 3 ; 0)$ và $B(5 ; 1 ;-2)$. Mặt phẳng trung trực của đoạn thẳng $A B$ có phương trình là\\
A. $2 x-y-z+5=0$.\\
B. $2 x-y-z-5=0$.\\
C. $x+y+2 z-3=0$.\\
D. $3 x+2 y-z-14=0$.

Câu 31. Họ tất cả các nguyên hàm của hàm số $f(x)=\frac{2 x-1}{(x+1)^{2}}$ trên khoảng $(-1 ;+\infty)$ là\\
A. $2 \ln (x+1)+\frac{2}{x+1}+C$.\\
B. $2 \ln (x+1)+\frac{3}{x+1}+C$.\\
C. $2 \ln (x+1)-\frac{2}{x+1}+C$.\\
D. $2 \ln (x+1)-\frac{3}{x+1}+C$.

Câu 32. Cho hàm số $f(x)$. Biết $f(0)=4$ và $f^{\prime}(x)=2 \cos ^{2} x+1, \forall x \in \mathbb{R}$, khi đó $\int_{0}^{\frac{\pi}{4}} f(x) d x$ bằng\\
A. $\frac{\pi^{2}+4}{16}$.\\
B. $\frac{\pi^{2}+14 \pi}{16}$.\\
C. $\frac{\pi^{2}+16 \pi+4}{16}$.\\
D. $\frac{\pi^{2}+16 \pi+16}{16}$.

Câu 33. Trong không gian $O x y z$, cho các điểm $A(1 ; 2 ; 0), B(2 ; 0 ; 2), C(2 ;-1 ; 3)$ và $D(1 ; 1 ; 3)$. Đường thẳng đi qua $C$ và vuông góc với mặt phẳng ( $A B D$ ) có phương trình là\\
A. $\left\{\begin{array}{l}x=-2-4 t \\ y=-2-3 t \\ z=2-t\end{array}\right.$.\\
B. $\left\{\begin{array}{l}x=2+4 t \\ y=-1+3 t \\ z=3-t\end{array}\right.$.\\
C. $\left\{\begin{array}{l}x=-2+4 t \\ y=-4+3 t \\ z=2+t\end{array}\right.$.\\
D. $\left\{\begin{array}{l}x=4+2 t \\ y=3-t \\ z=1+3 t\end{array}\right.$.

Câu 34. Cho số phức $z$ thỏa mãn $3(\bar{z}+i)-(2-i) z=3+10 i$. Mô đun của $z$ bằng\\
A. 3 .\\
B. 5 .\\
C. $\sqrt{5}$.\\
D. $\sqrt{3}$.

Câu 35. Cho hàm số $f(x)$, bảng xét dấu của $f^{\prime}(x)$ như sau:

\begin{center}
\begin{tabular}{c|ccccccccc}
$x$ & $-\infty$ &  & -3 &  & -1 &  & 1 &  & $+\infty$ \\
\hline
$f^{\prime}(x)$ &  & - & 0 & + & 0 & - & 0 & + &  \\
\hline
\end{tabular}
\end{center}

Hàm số $y=f(3-2 x)$ nghịch biến trên khoảng nào dưới đây?\\
A. $(4 ;+\infty)$.\\
B. (-2;1).\\
C. (2;4).\\
D. (1;2).

Câu 36. Cho hàm số $f(x)$, hàm số $y=f^{\prime}(x)$ liên tục trên $\mathbb{R}$ và có đồ thị như hình vẽ bên.\\
\includegraphics[max width=\textwidth, alt={}, center]{6a73def7-dd88-4569-925a-9b949b4179bd-4_460_641_1868_737}

Bất phương trình $f(x)<x+m$ ( $m$ là tham số thực) nghiệm đúng với mọi $x \in(0 ; 2)$ khi và chỉ khi\\
A. $m \geq f(2)-2$.\\
B. $m \geq f(0)$.\\
C. $m>f(2)-2$.\\
D. $m>f(0)$.

Câu 37. Chọn ngẫu nhiên 2 số tự nhiên khác nhau từ 25 số nguyên dương đầu tiên. Xác suất để chọn được hai số có tổng là một số chẵn bằng\\
A. $\frac{1}{2}$.\\
B. $\frac{13}{25}$.\\
C. $\frac{12}{25}$.\\
D. $\frac{313}{625}$.

Câu 38. Cho hình trụ có chiều cao bằng $5 \sqrt{3}$. Cắt hình trụ đã cho bởi mặt phẳng song song với trục và cách trục một khoảng bằng 1 , thiết diện thu được có diện tích bằng 30 . Diện tích xung quanh của hình trụ đã cho bằng\\
A. $10 \sqrt{3} \pi$.\\
B. $5 \sqrt{39} \pi$.\\
C. $20 \sqrt{3} \pi$.\\
D. $10 \sqrt{39} \pi$.

Câu 39. Cho phương trình $\log _{9} x^{2}-\log _{3}(3 x-1)=-\log _{3} m$ ( $m$ là tham số thực). Có tất cả bao nhiêu giá trị nguyên của $m$ để phương trình đã cho có nghiệm\\
A. 2 .\\
B. 4 .\\
C. 3 .\\
D. Vô số.

Câu 40. Cho hình chóp $S . A B C D$ có đáy là hình vuông cạnh $a$, mặt bên $S A B$ là tam giác đều và nằm trong mặt phẳng vuông góc với mặt phẳng đáy. Khoảng cách từ $A$ đến mặt phẳng ( $S B D$ ) bằng\\
A. $\frac{\sqrt{21} a}{14}$.\\
B. $\frac{\sqrt{21} a}{7}$.\\
C. $\frac{\sqrt{2} a}{2}$.\\
D. $\frac{\sqrt{21} a}{28}$.

Câu 41. Cho hàm số $f(x)$ có đạo hàm liên tục trên $\mathbb{R}$. Biết $f(4)=1$ và $\int_{0}^{1} x f(4 x) \mathrm{d} x=1$, khi đó $\int_{0}^{4} x^{2} f^{\prime}(x) \mathrm{d} x$ bằng\\
A. $\frac{31}{2}$.\\
B. -16 .\\
C. 8 .\\
D. 14 .

Câu 42. Trong không gian $O x y z$, cho điểm $A(0 ; 4 ;-3)$. Xét đường thẳng $d$ thay đổi, song song với trục $O z$ và cách trục $O z$ một khoảng bằng 3. Khi khoảng cách từ $A$ đến $d$ nhỏ nhất, $d$ đi qua điểm nào dưới đây?\\
A. $P(-3 ; 0 ;-3)$.\\
B. $M(0 ;-3 ;-5)$.\\
C. $N(0 ; 3 ;-5)$.\\
D. $Q(0 ; 5 ;-3)$.

Câu 43. Cho hàm số bậc ba $y=f(x)$ có đồ thị như hình vẽ bên.\\
\includegraphics[max width=\textwidth, alt={}, center]{6a73def7-dd88-4569-925a-9b949b4179bd-5_357_471_1062_822}

Số nghiệm thực của phương trình $\left|f\left(x^{3}-3 x\right)\right|=\frac{4}{3}$ là\\
A. 3 .\\
B. 8 .\\
C. 7 .\\
D. 4 .

Câu 44. Xét các số phức $z$ thỏa mãn $|z|=\sqrt{2}$. Trên mặt phẳng tọa độ $O x y$, tập hợp điểm biểu diễn của các số phức $\mathrm{w}=\frac{4+i \mathrm{z}}{1+z}$ là một đường tròn có bán kính bằng\\
A. $\sqrt{34}$.\\
B. 26 .\\
C. 34 .\\
D. $\sqrt{26}$.

Câu 45. Cho đường thẳng $y=x$ và Parabol $y=\frac{1}{2} x^{2}+a$ ( $a$ là tham số thực dương). Gọi $S_{1}$ và $S_{2}$ lần lượt là diện tích của hai hình phẳng được gạch chéo trong hình vẽ bên. Khi $S_{1}=S_{2}$ thì $a$ thuộc khoảng nào sau đây?\\
\includegraphics[max width=\textwidth, alt={}, center]{6a73def7-dd88-4569-925a-9b949b4179bd-5_460_510_2037_804}\\
A. $\left(\frac{3}{7} ; \frac{1}{2}\right)$.\\
B. $\left(0 ; \frac{1}{3}\right)$.\\
C. $\left(\frac{1}{3} ; \frac{2}{5}\right)$.\\
D. $\left(\frac{2}{5} ; \frac{3}{7}\right)$.

Câu 46. Cho hàm số $f(x)$, bảng biến thiên của hàm số $f^{\prime}(x)$ như sau

\begin{center}
\begin{tabular}{l|lllll}
$x$ & $-\infty$ & -1 & 0 & 1 & $+\infty$ \\
\hline
$y$ & $+\infty$ &  & ${ }^{2}$ &  & $+\infty$ \\
\hline
\end{tabular}
\end{center}

Số điểm cực trị của hàm số $y=f\left(x^{2}-2 x\right)$ là\\
A. 9 .\\
B. 3 .\\
C. 7 .\\
D. 5 .

Câu 47. Cho lăng trụ $A B C \cdot A^{\prime} B^{\prime} C^{\prime}$ có chiều cao bằng 8 và đáy là tam giác đều cạnh bằng 6. Gọi $M, N$ và $P$ lần lượt là tâm của các mặt bên $A B B^{\prime} A^{\prime}, A C^{\prime} A^{\prime}$ và $B C C^{\prime} B^{\prime}$. Thể tích của khối đa diện lồi có các đỉnh là các điểm $A, B, C, M, N, P$ bằng:\\
A. $27 \sqrt{3}$.\\
B. $21 \sqrt{3}$.\\
C. $30 \sqrt{3}$.\\
D. $36 \sqrt{3}$.

Câu 48. Trong không gian $O x y z$, cho mặt cầu $(S): x^{2}+y^{2}+(z+\sqrt{2})^{2}=3$. Có tất cả bao nhiêu điểm $A(a ; b ; c)(a, b, c$ là các số nguyên) thuộc mặt phẳng ( $O x y$ ) sao cho có ít nhất hai tiếp tuyến của ( $S$ ) đi qua $A$ và hai tiếp tuyến đó vuông góc với nhau?\\
A. 12 .\\
B. 8 .\\
C. 16 .\\
D. 4 .

Câu 49. Cho hai hàm số $y=\frac{x-3}{x-2}+\frac{x-2}{x-1}+\frac{x-1}{x}+\frac{x}{x+1}$ và $y=|x+2|-x+m$ ( $m$ là tham số thực) có đồ thị lần lượt là $\left(C_{1}\right)$ và $\left(C_{2}\right)$. Tập hợp tất cả các giá trị của $m$ để $\left(C_{1}\right)$ và $\left(C_{2}\right)$ cắt nhau tại 4 điểm phân biệt là\\
A. $(-\infty ; 2]$.\\
B. $[2 ;+\infty)$.\\
C. $(-\infty ; 2)$.\\
D. $(2 ;+\infty)$.

Câu 50. Cho phương trình $\left(4 \log _{2}^{2} x+\log _{2} x-5\right) \sqrt{7^{x}-m}=0$ ( $m$ là tham số thực). Có tất cả bao nhiêu giá trị nguyên dương của $m$ để phương trình đã cho có đúng hai nghiệm phân biệt\\
A. 49 .\\
B. 47 .\\
C. Vô số.\\
D. 48.


\end{document}